\section{The shortcut}
\label{sec:shortcut}
\textbf{Source is predictable from text.} We train a word TF-IDF + logistic regression classifier to
predict a skill's \emph{source benchmark} (not its label) and evaluate it with held-out skills. It
reaches \ProbeAcc{} accuracy over \ProbeSources{} sources (chance $1/\ProbeSources{}$). Since most
sources carry a single label, a malicious/benign classifier can score well by recognising the source
and never learning what makes a skill dangerous.

\textbf{A concrete shortcut feature.} The token \texttt{openclaw}, platform boilerplate from skills
published for the OpenClaw agent, is one of the strongest ``benign'' cues an all-data model learns.
Yet inside the one source that contains both labels at scale (SkillTrustBench) it appears about
equally in benign and malicious skills, so it carries no signal about maliciousness. It looks benign
only because benign training data is dominated by sources where the token is common, while the
largest malicious source rarely uses it. The cue is an artifact of how the benchmark's sources were
mixed, exactly the mechanism of \cite{arp2022,torralba2011}.

\section{Training where both labels share a source}
\label{sec:fix}
If a source contains only one label it can teach only style; if it contains both, style cannot
separate the labels inside it, so the model must use content. We therefore train only on the sources
that contain both labels. On the unseen-source test set this cuts the benign false-positive rate from
\FprAll{} to \FprMixed{} (95\% CI \FprMixedCI{}) and raises PR-AUC from \PrAll{} to \PrMixed{}
(Table~\ref{tab:fix}).

\begin{table}[t]
\caption{Word TF-IDF + logistic regression on the unseen-source test set, default $0.5$ threshold.}
\label{tab:fix}
\centering
\begin{tabular}{lcc}
\toprule
training data & benign FPR & PR-AUC \\
\midrule
all sources & \FprAll{} & \PrAll{} \\
size-matched random control & \FprSizeMatched{} & -- \\
mixed-label sources only & \FprMixed{} & \PrMixed{} \\
\bottomrule
\end{tabular}
\end{table}

\textbf{It is not simply ``less data''.} Mixed-only training uses far fewer skills than all-data. A
control that draws a random subset of all-data of the same size and label balance still flags
\FprSizeMatched{} of unseen-source benign skills, so the gain comes from \emph{which} data, not how
much. Re-weighting every (source, label) group equally does not help either: the style signal is
still present, only reweighted.

\textbf{Thresholds do not transfer.} A threshold tuned to flag 2\% of \emph{validation} benign skills
(validation comes from the training sources) flags far more on the unseen test sources, so any number
reported ``at $k\%$ FPR'' must say where the threshold was set. We report both an oracle operating
point (threshold on the test set, an upper bound) and the deployable one (threshold on validation).

\section{What the fix learns}
\label{sec:mechanism}
Mixed-only training is dominated by one source, SkillTrustBench (\MixedSkillTrustBench{} of the
mixed-label training skills; AgentTrap contributes \MixedAgentTrap{} and SkillTrojan
\MixedSkillTrojan{}), and that source builds most malicious skills by injecting a payload into a real
benign skill. So the fix could work for two different reasons: label mixing within a source, or the
presence of minimal pairs (a benign skill and a near-copy carrying a payload). We separate them.

\MechanismSentence{}

\textbf{Which sources bring the false positives back.} Adding one single-label source at a time to
the mixed-only training set locates the data responsible for the collapse: the largest single-label
malicious sources restore the high benign false-positive rate, while small ones do not. This shows
the failure is driven by the single-label malicious bulk, not by a lack of benign coverage alone.

\textbf{Masking the obvious token is not enough.} Removing the \texttt{openclaw} token, or stripping
the YAML frontmatter, from all-data training does not close the gap: the model falls back on the many
other style cues that distinguish the sources. The shortcut is distributed across the text, so
editing one feature cannot fix an all-data model; changing the training composition can.

\section{Fine-tuned transformer}
\label{sec:transformer}
To check that the shortcut is a property of the data and not of linear models, we fine-tune
ModernBERT-base on both training sets and score it on the same unseen-source test skills, comparing
every model with a paired, class-stratified bootstrap. \TransformerSentence{} Per-source rates on the
small benign test groups (ATR, SkillFortifyBench) are reported with exact Clopper-Pearson intervals,
since ATR contributes only \TestAtrBenign{} of the \TestBenign{} benign test skills.
